% ============================================================
% ch13.tex —— 第 13 章「非退化谱网络」
% 印刷 pp.141–162（PDF p181–p202，22 页）
% 底本：Springer LNM 2386, Ch.13 Non-degenerate spectral networks
% 本块无编号公式（tags2.txt 中该页段为空，已逐页核对图像）
% ============================================================

\chapter{非退化谱网络}
% [原文 p.141]

\ChapterAbstract{谱网络的引入以非阿贝尔化映射的构造为动机。随后我们定义纤毛曲面上的谱网络。我们从最简单的小谱网络情形开始，然后给出非退化谱网络更一般的定义。我们给出谱网络构造的两种不同进路：第一种是纯拓扑与组合的，适用于带穿孔、且不带任何附加结构的纤毛曲面 $S$。第二种称为 WKB 谱网络，它更为解析，需要在 $S$ 上选定一个复结构。并非所有一般谱网络都能用这种解析构造实现。此外，我们描述一条用谱网络作道路提升的规则，它推广了非分歧覆叠通常的道路提升性质，并且对分歧覆叠而言是同伦不变的。}

\section{迈向非阿贝尔化映射}

设 $S=S_{g,p}$ 为纤毛曲面。考虑基本群表示的空间
\[
\mathcal{E}(S,\GL_n(\C)):=\Hom\bigl(\pi_1(S),\GL_n(\C)\bigr)/\GL_n(\C),
\]
其中 $\GL_n(\C)$ 通过共轭作用在同态空间上。

\begin{example}[13.1]
在最简单而非平凡的情形 $n=1$，群 $\GL_1(\C)=\C^\times$ 阿贝尔，$\GL_1(\C)$ 通过共轭在同态空间上的作用平凡。我们得到
\[
\mathcal{E}(S,\GL_1(\C)):=\Hom\bigl(\pi_1(S),\C^\times\bigr).
\]
由于 $\C^\times$ 阿贝尔，任何同态 $\rho\in\Hom(\pi_1(S),\C^\times)$ 把 $\pi_1(S)$ 中所有换位子都送到单位元，即它穿过第一同调群 $H_1(S,\Z)=\pi_1(S)/[\pi_1(S),\pi_1(S)]$：
% [原文 p.142]
\[
\Hom\bigl(\pi_1(S),\C^\times\bigr)=\Hom\bigl(H_1(S,\Z),\C^\times\bigr).
\]
\end{example}

由于 $H_1(S,\Z)$ 是 $\Z$ 上的自由模，存在 $H_1(S,\Z)$ 的一组基 $\mathbf{e}:=(e_1,\dots,e_s)$。于是映射
\[
\begin{array}{rcl}
\Hom\bigl(H_1(S,\Z),\C^\times\bigr)&\to&(\C^\times)^s\\
\rho&\mapsto&(\rho(e_1),\dots,\rho(e_s))
\end{array}
\]
给出这两个空间之间的同胚。

% [Tutor Note: 维数核对的细节：若 $\Sigma\to S$ 是非分歧 $n:1$ 覆叠，由 Riemann–Hurwitz 有 $2g_\Sigma-2=n(2g-2)$；又 $\dim\mathcal{E}(\Sigma,\C^\times)=2g_\Sigma$（因 $H_1(\Sigma,\Z)$ 的秩为 $2g_\Sigma$），故 $2g_\Sigma=n(2g-2)+2$。原文写作 $n(2g-2)$，与精确值相差常数 $2$；这不影响「前者远小于 $n^2(2g-2)$」的结论，但引用维数时宜按 $2g_\Sigma$ 核对。]

上例说明：若 $G=\C^\times$（或更一般地，$G$ 为任何阿贝尔群），则空间 $\mathcal{E}(S,G)$ 的拓扑容易理解。为了研究 $G$ 非阿贝尔的一般情形，我们希望把它化归到阿贝尔情形。更确切地说，我们要找另一个曲面 $\Sigma$ 以及一个映射
\[
\mathcal{E}(\Sigma,\C^\times)\to\mathcal{E}(S,\GL_n(\C)),
\]
它“足够好”以理解空间 $\mathcal{E}(S,\GL_n(\C))$，例如是一个一般地有限对一、像为开稠密集的映射。我们把任何这样的映射称为\emph{非阿贝尔化映射}（non-abelianization map）。我们也常考虑与这些非阿贝尔化映射局部互逆的映射。这类映射称为\emph{阿贝尔化映射}（abelianization maps）
\[
\mathcal{E}(S,\GL_n(\C))\to\mathcal{E}(\Sigma,\C^\times).
\]

我们来看一个倾向于成为非阿贝尔化映射、但按上述意义并不“足够好”的映射的例子。

\begin{example}[13.2]
如我们在定理 3.30 中所见，空间 $\mathcal{E}(S,\GL_n(\C))$ 可与秩 $n$ 的平坦 $\C$-向量丛的等价类集合等同。设 $\pi\colon\Sigma\to S$ 为一个（非分歧的）$n:1$-覆叠映射。该映射诱导一个映射
\[
\begin{array}{rcl}
\pi_*\colon\mathcal{E}(\Sigma,\C^\times)&\to&\mathcal{E}(S,\GL_n(\C))\\
(L,\nabla)&\mapsto&(\pi_*L,\pi_*\nabla)
\end{array}
\]
其中 $L$ 是 $\Sigma$ 上带有平坦联络 $\nabla$ 的秩 $1$ 的 $\C$-丛，而 $\pi_*L$ 是 $S$ 上秩 $n$ 的丛，配有我们将在下文描述的平坦联络 $\pi_*\nabla$。

对点 $z\in S$，令 $\pi^{-1}(z)=\{z^{(1)},\dots,z^{(n)}\}$ 为其原像。我们定义丛 $\pi_*L$ 在点 $z\in S$ 上的纤维为
\[
(\pi_*L)_z:=\bigoplus_{i=1}^{n}L_{z^{(i)}},
\]
其中 $L_{z^{(i)}}$ 是 $L$ 在 $z^{(i)}\in\Sigma$ 上的纤维。

现在令 $\gamma\colon[0,1]\to S$ 为 $S$ 中起点为 $z$ 的道路，即 $\gamma(0)=z$。要定义 $S$ 上的联络，我们需要描述沿 $\gamma$ 的平行移动。为此，我们把 $\gamma$ 提升到 $\Sigma$。恰有 $n$ 条提升 $\gamma^{(1)},\dots,\gamma^{(n)}$，使得对每个 $i\in\{1,\dots,n\}$ 都有 $\gamma^{(i)}(0)=z^{(i)}$。
\end{example}

% [原文 p.143]
沿每条 $\gamma^{(i)}$ 的平行移动由 $\nabla$ 定义，于是我们得到 $n$ 个线性映射 $P_{\gamma^{(i)}}\colon L_{z^{(i)}}\to L_{\gamma^{(i)}(1)}$。我们如下定义沿 $\gamma$ 的平行移动：
\[
P_\gamma:=\bigoplus_{i=1}^{n}P_{\gamma^{(i)}}\colon
\bigoplus_{i=1}^{n}L_{z^{(i)}}\longrightarrow\bigoplus_{i=1}^{n}L_{\gamma^{(i)}(1)},
\qquad
(v_1,\dots,v_n)\mapsto\bigl(P_{\gamma^{(1)}}v_1,\dots,P_{\gamma^{(n)}}v_n\bigr).
\]
该映射定义了 $\pi_*L$ 上的联络 $\pi_*\nabla$，它自动平坦，因为 $\nabla$ 平坦（习题）。

映射 $\pi_*$，若更一般地考虑不一定约化的基本群表示，可能是非阿贝尔化映射的好候选；但更详细的计算表明，空间 $\mathcal{E}(\Sigma,\C^\times)$ 的复维数是 $n(2g-2)$（其中 $g$ 是 $S$ 的亏格），而 $\mathcal{E}(S,\GL_n(\C))$ 的维数是 $n^2(2g-2)$。所以 $\pi_*$ 没有机会有开的像。然而这个例子给出了如何构造非阿贝尔化映射的想法。在下一节我们将看到，若用分歧覆叠代替非分歧覆叠，非阿贝尔化映射确实可以构造出来。这将引出谱网络的概念。

\section{路径代数}

如前节所宣布，我们将用分歧覆叠来构造非阿贝尔化映射。分歧覆叠的问题在于：对它们而言，标准的道路提升规则不是同伦不变的。我们将在下一个例子中说明这一问题：

\begin{example}[13.3]
回顾注 8.13 中给出的例子，它说明标准的道路提升规则对分歧覆叠不是同伦不变的。这个例子也说明：若我们像例 13.2 那样，试图把平坦联络从分歧覆叠 $\Sigma$ 前推到 $S$ 中分支轨迹 $B\subset S$ 之外，则所得到的 $S\setminus B$ 上的联络不会延拓为 $S$ 上的平坦联络，因为它在分支点周围会有非平凡的单值性。
\end{example}

我们将修改例 13.2 的构造，使之适用于分歧覆叠。为此，我们需要曲面上的路径代数这一概念。

\begin{definition}[13.4]
曲面 $S$ 上的\emph{道路空间}定义为
\[
\mathcal{P}(S):=\bigl\{\gamma\colon[0,1]\to S\ \text{连续}\bigr\}/\sim,
\]
其中两条道路 $\gamma_0,\gamma_1$ 等价（记作 $\gamma_0\sim\gamma_1$），如果 $\gamma_0$ 与 $\gamma_1$ 相对于端点同伦，即 $\gamma_0(0)=\gamma_1(0)$、$\gamma_0(1)=\gamma_1(1)$，且存在连续映射 $H\colon[0,1]\times[0,1]\to S$ 使得
\begin{itemize}
\item $H(0,t)=\gamma_0(t)$，$H(1,t)=\gamma_1(t)$ 对一切 $t\in[0,1]$，
% [原文 p.144]
\item $H(s,0)=\gamma_0(0)=\gamma_1(0)$，$H(s,1)=\gamma_0(1)=\gamma_1(1)$ 对一切 $s\in[0,1]$。
\end{itemize}
\end{definition}

\begin{definition}[13.5]
$S$ 上的\emph{路径代数}是 $\mathcal{P}(S)$ 上的自由阿贝尔群 $\bigl(\Z[\mathcal{P}(S)],+\bigr)$，赋予如下定义的乘法 $\cdot$（图 13.1）：若 $\gamma_1(1)\neq\gamma_2(0)$，则 $\gamma_2\cdot\gamma_1=0$；否则 $\gamma_2\cdot\gamma_1$ 为 $\gamma_1$ 与 $\gamma_2$ 的拼接，即 $\gamma_2\cdot\gamma_1=\gamma$，其中
\[
\gamma(t)=
\begin{cases}
\gamma_1(2t), & \text{当 } t\in\bigl[0,\tfrac12\bigr];\\[3pt]
\gamma_2(2t-1), & \text{当 } t\in\bigl[\tfrac12,1\bigr].
\end{cases}
\]
\end{definition}

\Fig[0.72\linewidth]{figures/fig_13_1.png}{13.1}{路径代数中的乘法。}

集合 $\mathcal{P}(S)$ 自然地嵌入 $\Z[\mathcal{P}(S)]$。

设 $\pi\colon\Sigma\to S$ 为分歧覆叠。我们用 $B'$ 表示 $\Sigma$ 中分歧点的集合，并令 $B:=\pi(B')$。设 $b\in B'$，$D$ 是 $b$ 的一个不含其他分歧点的开可缩邻域，且 $z\in D\setminus\{b\}$。令 $\alpha_z^b\colon[0,1]\to D\setminus\{b\}$ 为 $z$ 处的一个环路，它绕 $b$ 一周且在 $D$ 中可缩。环路 $\alpha_z^b$ 作为道路并不唯一，但所有这样的 $\alpha_z^b$ 都同伦，即所有这样的 $\alpha_z^b$ 定义了 $\mathcal{P}(\Sigma\setminus B')$ 中同一个元素，我们仍（略微滥用记号地）把它记作 $\alpha_z^b$。此外，令 $e_z\colon[0,1]\to D\setminus\{b\}$ 为 $z$ 处的一个在 $D\setminus\{b\}$ 中可缩的环路。类似地，$e_z$ 作为道路并不唯一，但所有这样的 $e_z$ 都同伦，即都定义 $\mathcal{P}(\Sigma\setminus B')$ 中同一个元素，我们仍把它记作 $e_z$。

元素 $e_z$ 可以解释为“点 $z$ 处的单位元”，即对每个满足 $\gamma(0)=z$ 的元素 $\gamma\in\mathcal{P}(\Sigma\setminus B')$，有 $\gamma\cdot e_z=\gamma$。

我们定义由一切 $\alpha_z^b+e_z$（$z\in\Sigma$，$b\in B'$）生成的理想 $\mathcal{I}$。

\begin{definition}[13.6]
\emph{扭转路径代数}是商代数
\[
\Z^{\mathrm{tw}}_{B'}\bigl[\mathcal{P}(\Sigma)\bigr]=\Z\bigl[\mathcal{P}(\Sigma\setminus B')\bigr]/\mathcal{I}.
\]
\end{definition}

我们将把扭转路径代数 $\Z^{\mathrm{tw}}_{B'}[\mathcal{P}(\Sigma)]$ 的元素写成 $\Z[\mathcal{P}(\Sigma\setminus B')]$ 的元素，同时记住需要等同某些元素。例如在 $\Z^{\mathrm{tw}}_{B'}[\mathcal{P}(\Sigma)]$ 中有等式 $\alpha_z^b=-e_z$。

% [原文 p.145]
\section{小谱网络}

设 $\bar S$ 为带有有限标记点集 $P=\{s_1,\dots,s_n\}$（$n\in\Z_{>0}$）的紧曲面。令 $S=\bar S\setminus P$ 为相应的纤毛曲面。

设 $\pi\colon\bar\Sigma\to\bar S$ 为 $\bar S$ 的分歧 $k$ 重覆叠。我们也考虑相应的限制 $\pi\colon\Sigma\to S$，其中 $\Sigma=\pi^{-1}(S)=\bar\Sigma\setminus P'$，且 $P'=\pi^{-1}(P)$ 是 $\Sigma$ 的穿孔集。我们假定覆叠在穿孔上不分歧。如前，用 $B'$ 表示 $\Sigma$ 的分歧点集，$B=\pi(B')\subset S$。

我们假定每个分支点都是\emph{简单的}。这意味着每个 $b\in B$ 有开邻域 $U$，使得
\[
\pi^{-1}(U)=\bigsqcup_{i=1}^{k-1}V_i,
\]
其中 $V_i$ 是 $b$ 在 $\Sigma$ 中某个原像的开邻域，$\pi|_{V_i}\colon V_i\to U$ 对 $i=1,\dots,k-2$ 是同胚，而映射 $\pi|_{V_{k-1}}\colon V_{k-1}\to U$ 在同痕意义下是映射 $z\mapsto z^2$。

我们还假定：对每个 $s\in P$，$\pi^{-1}(s)=\{s_i^{(1)},\dots,s_i^{(k)}\}$ 上给定了一个次序，由 $s_i^{(1)}<s_i^{(2)}<\dots<s_i^{(k)}$ 给出。纤维上的这个次序是覆叠上的一个附加结构。

现在准备定义小谱网络，它是谱网络的一个简化但仍有用的版本。我们采用 [12] 中的定义，它与 Gaiotto--Moore--Neitzke 在 [154] 中的原始定义等价。

\begin{definition}[13.7]
与覆叠 $\pi$ 相关联的秩 $k$ 的\emph{小谱网络} $\mathcal{W}$ 是 $\bar\Sigma$ 上的一个图，即一有限族光滑、简单、定向的道路 $\{p_h\colon[-1,1]\to\bar\Sigma\}_{h\in H}$，其中 $H$ 是有限指标集（图 13.2），这些道路只能横截相交，并满足以下条件：
\begin{enumerate}
\item $\pi(p_h(t))=\pi(p_h(-t))$ 对一切 $t\in[-1,1]$；
\item $p_h(0)$ 是 $\Sigma$ 上的分歧点；
\item $\pi(p_h(1))=\pi(p_h(-1))=s_j\in P$（对某个 $j$）是穿孔，且相对于 $\pi^{-1}(s_j)$ 上的次序有 $p_h(-1)<p_h(1)$；
% [原文 p.146]
\item 每个分支点在 $S$ 中、每个分歧点在 $\Sigma$ 中都有如图 13.3 的邻域。
\item 不同的道路 $p_h$ 与 $p_{h'}$ 只能在分歧点处、穿孔处相交，或如图 13.4 那样相交。在后一情形，当交点为 $p_h(t)=p_{h'}(t')$（$t,t'\in(-1,1)\setminus\{0\}$）时，我们要求 $\mathcal{W}$ 中没有线经过对径点 $p_h(-t)$ 或 $p_{h'}(-t')$。
\end{enumerate}
\end{definition}

\Fig[0.55\linewidth]{figures/fig_13_2.png}{13.2}{谱网络的一条线：$\bar S$ 中（左）与 $\bar\Sigma$ 中（右）。}

\Fig[0.50\linewidth]{figures/fig_13_3.png}{13.3}{$S$ 中分支点（左）与 $\Sigma$ 中分歧点（右）的邻域。}

\Fig[0.75\linewidth]{figures/fig_13_4.png}{13.4}{谱网络诸线的相交。图为 $\bar\Sigma$ 中的情形。}

我们记
\[
\mathcal{W}_\Sigma:=\bigcup_{h\in H}p_h\bigl((-1,1)\bigr)\subseteq\Sigma
\]
为 $\Sigma$ 上的谱网络，并记
\[
\mathcal{W}_S:=\pi(\mathcal{W}_\Sigma)\subseteq S
\]
为 $S$ 上的谱网络。

\begin{remark}[13.8]
当不会引起歧义时，我们用 $\mathcal{W}$ 表示 $S$ 上或 $\Sigma$ 上的谱网络。
\end{remark}

\begin{remark}[13.9]
对小谱网络，如图 13.5 中的道路相交（称为\emph{关节}（joint））是被禁止的。我们将在 §13.4 定义一类更一般的谱网络，其中允许这种类型的相交。
\end{remark}

% [原文 p.147]
\Fig[0.45\linewidth]{figures/fig_13_5.png}{13.5}{被禁止的道路相交。图为 $\bar\Sigma$ 中的情形。}

\begin{proposition}[13.10]
注 13.9 与定义 13.7 (5) 中所述的相交，在秩 $2$ 的谱网络中不可能出现。
\end{proposition}

\begin{proof}
这一事实直接由定义 13.7 中的公理 (5) 得出。若假定存在两条道路 $p_h$ 与 $p_{h'}$，使得存在 $t,t'\in(0,1)\setminus\{0\}$ 且
\[
\pi(p_h(t))=\pi(p_h(-t))=\pi(p_{h'}(t'))=\pi(p_{h'}(t')),
\]
则因为 $\pi\colon\Sigma\to S$ 恰有两叶，只有两种可能：
\[
p_h(t)=p_{h'}(t')\quad\text{且}\quad p_h(-t)=p_{h'}(-t')
\]
或
\[
p_h(t)=p_{h'}(-t')\quad\text{且}\quad p_h(-t)=p_{h'}(t'),
\]
见图 13.6 与图 13.7。这两种情形都被定义 13.7 的公理 (5) 禁止。
\end{proof}

\Fig[0.50\linewidth]{figures/fig_13_6.png}{13.6}{图为 $S$ 中的情形。}

% [Tutor Note: 命题 13.10 证明中的等式 $\pi(p_h(t))=\pi(p_h(-t))=\pi(p_{h'}(t'))=\pi(p_{h'}(t'))$ 末项重复了 $t'$，按上下文（交点满足 $p_h(t)=p_{h'}(t')$，并需要关于对径点的信息）应为 $\pi(p_{h'}(-t'))$；紧接着给出的两种可能与图 13.6、13.7 也都用到 $p_{h'}(-t')$。]

\begin{remark}[13.11]
由命题 13.10 可知，所有秩 $2$ 的谱网络都是小的。
\end{remark}

% [原文 p.148]
\Fig[0.42\linewidth]{figures/fig_13_7.png}{13.7}{图为 $\Sigma$ 中的情形。}

\section{非退化谱网络}

正如我们在注 13.9 中所见，谱网络的线之间存在称为关节的相交，这类相交对小谱网络是被禁止的。为了处理这类相交，我们需要在每次出现这种相交时向谱网络添加额外的道路。

如前，设 $\pi\colon\bar\Sigma\to\bar S$ 为分歧 $k$ 重覆叠。我们假定 $k\ge3$，因为如我们所见，秩 $2$ 的谱网络中不会出现关节。

\begin{definition}[13.12]
与 $\pi$ 相关联的秩 $k$ 的\emph{非退化谱网络} $\mathcal{W}$ 由 $\Sigma$ 上两族有限的光滑、简单、定向的道路组成：$\{p_h\colon[-1,1]\to\bar\Sigma\}_{h\in H_1}$ 与 $\{p_h\colon[-1,1]\setminus\{0\}\to\bar\Sigma\}_{h\in H_2}$，其中 $H_1$ 与 $H_2$ 是两个不相交的指标集。这些道路只能横截相交，并满足以下条件：
\begin{enumerate}
\item 对一切 $h\in H_1\cup H_2$ 与一切 $t\in(0,1]$，$\pi(p_h(t))=\pi(p_h(-t))$；
\item 对一切 $h\in H_1$，$\pi(p_h(0))\in B$ 是分支点；
\item 对一切 $h\in H_2$，存在极限 $p_h(+0):=\lim_{t\to+0}p_h(t)$ 与 $p_h(-0):=\lim_{t\to-0}p_h(t)$。此外 $p_h(+0)\neq p_h(-0)$，但 $\pi(p_h(+0))=\pi(p_h(-0))$；
\item 对一切 $h\in H_1\cup H_2$，$\pi(p_h(1))=\pi(p_h(-1))=s_j\in P$（对某个 $j$）是穿孔，且相对于 $\pi^{-1}(s_j)$ 上的次序有 $p_h(-1)<p_h(1)$；
\item 任何分歧点的邻域如图 13.8。
% [原文 p.149]
\item $\bar\Sigma$ 中不同的道路 $p_h$ 与 $p_{h'}$ 只能在一个点 $p_h(t)=p_{h'}(t')$ 处相遇，且下列之一成立：
\begin{enumerate}
\item $h,h'\in H_1$，$t=t'=0$，且 $p_h(t)=p_{h'}(t')$ 是分歧点；
\item $h,h'\in H_1\cup H_2$，$t,t'\in\{-1,1\}$，且 $p_h(t)=p_{h'}(t')$ 是穿孔；
\item $h,h'\in H_1\cup H_2$，$t\cdot t'>0$，此时我们还要求 $\mathcal{W}$ 中没有线经过 $p_h(-t)$ 或 $p_{h'}(-t')$；
\item $h,h'\in H_1\cup H_2$，$t\cdot t'<0$，此时不妨设 $t<0$、$t'>0$。在这种情形下，存在道路 $p_{h''}$（$h''\in H_2$）使得 $p_{h''}(+0)=p_h(-t)$ 且 $p_{h''}(-0)=p_{h'}(-t')$。此外，$\mathcal{W}$ 中没有其他线经过 $p_h(-t)$ 或 $p_{h'}(-t')$。这种类型的相交称为\emph{关节}。
\end{enumerate}
\end{enumerate}
\end{definition}

% [Tutor Note: 条件 (3) 说明了 $H_2$ 型道路为何定义在 $[-1,1]\setminus\{0\}$ 上：它们在 $t=0$ 处不连续，左右极限落在 $\Sigma$ 的不同叶中（但投影到 $S$ 的同一点），这正是关节的模样。作图核对时注意 $\pi(p_h(+0))=\pi(p_h(-0))$ 而 $p_h(+0)\ne p_h(-0)$。]

\Fig[0.50\linewidth]{figures/fig_13_8.png}{13.8}{$S$ 中分支点（左）与 $\Sigma$ 中分歧点（右）的邻域。}

我们记
\[
\mathcal{W}_\Sigma:=\bigcup_{h\in H_1\cup H_2}p_h\bigl((-1,1)\bigr)\subseteq\Sigma
\]
为 $\Sigma$ 上的谱网络，并记
\[
\mathcal{W}_S:=\pi(\mathcal{W}_\Sigma)\subseteq S
\]
为 $S$ 上的谱网络。

\begin{example}[13.13]
图 13.9 中给出的例子是一个非退化但非小的谱网络，因为它含有两个关节。
\end{example}

\Fig[0.42\linewidth]{figures/fig_13_9.png}{13.9}{一个非退化但非小的谱网络。}

% [原文 p.150]
\begin{remark}[13.14]
小谱网络特别地也是非退化的。
\end{remark}

\begin{remark}[13.15]
与小谱网络相反，非退化谱网络不必是有限的。鉴于本书的范围，我们将总假定 $S$ 的任何紧子集只与谱网络相交有限多次。若 $H_1$ 与 $H_2$ 是有限集，我们称相应的谱网络是\emph{有限的}。
\end{remark}

\section{组合构造}

本小节详细给出谱曲面与某个小谱网络的组合构造，后者与底曲面的一个三角剖分相关联；该构造先在秩 $2$ 情形演示，然后推广到高阶。

\subsection{秩二谱网络}

设 $S$ 为纤毛曲面，$\Delta$ 为 $S$ 的一个理想三角剖分。$\Delta$ 在 $S$ 中的补集是一族可缩开集，它们恰与三个穿孔相接。我们称这些集合在 $S$ 中的闭包为 $\Delta$ 的\emph{三角形}。用 $\Delta_3$ 表示 $\Delta$ 的所有三角形构成的集合。注意三角形是 $S$ 的闭可缩子集，且每个三角形的边界由 $\Delta$ 的三条边组成。

对每个三角形 $T\in\Delta_3$，我们固定一个同胚 $\phi_T\colon T\to\mathbb{T}$，其中 $\mathbb{T}$ 是 $\C$ 中挖去顶点 $1$、$\ee^{2i\pi/3}$ 与 $\ee^{4i\pi/3}$ 的紧正则欧氏三角形。我们称 $\mathbb{T}$ 为\emph{标准三角形}。

现在令 $\HH$ 为 $\C$ 中挖去顶点 $1$、$\ee^{i\pi/3}$、$\ee^{2i\pi/3}$、$-1$、$\ee^{4i\pi/3}$ 与 $\ee^{5i\pi/3}$ 的六边形。$\C$ 上的映射 $z\mapsto z^2$ 把 $\HH$ 映到 $\mathbb{T}$（见图 13.10）。该映射是只有一个分歧点 $z=0$ 的分歧覆叠。现在对每个 $T\in\Delta_3$，考虑对 $H_T:=(\HH,T)$，并定义映射 $\psi_T\colon H_T\to\HH$ 为 $(z,T)\mapsto z$。于是我们平凡地得到如下交换图表，以及映射 $\theta\colon H_T\to T:=\psi_T\circ(z\mapsto z^2)\circ\phi_T^{-1}$，它是分歧覆叠：
\begin{center}
\begin{tikzpicture}[scale=1.25]
  \node (HT) at (0,1.25)  {$H_T$};
  \node (H)  at (2.9,1.25) {$\HH$};
  \node (T)  at (0,0)     {$T$};
  \node (TT) at (2.9,0)   {$\mathbb{T}$};
  \draw[->] (HT) -- node[above] {$\psi_T$} (H);
  \draw[->] (T)  -- node[below] {$\phi_T$} (TT);
  \draw[->] (HT) -- node[left]  {$\theta_T$} (T);
  \draw[->] (H)  -- node[right] {$z\mapsto z^2$} (TT);
\end{tikzpicture}
\end{center}

在 $\HH$ 中，我们取三条定向道路：
% [原文 p.151]
\[
\begin{array}{rcl}
p_k\colon(-1,1)&\to&\HH\\
t&\mapsto&t\,\ee^{2ik\pi/3},
\end{array}
\]
对 $k\in\{0,1,2\}$，并用映射 $z\mapsto z^2$ 投影每条道路。显然，这些道路满足覆叠 $\HH\to\mathbb{T}$ 的谱网络条件。

\Fig[0.78\linewidth]{figures/fig_13_10.png}{13.10}{覆叠 $\HH\to\mathbb{T}$。}

我们希望按下述方式沿边粘合六边形 $H_T$ 得到 $\Sigma$：设 $e$ 是 $\Delta$ 的一条边，它分隔两个三角形 $T_1$ 与 $T_2$。令 $\alpha_e\colon(0,1)\to e$ 为 $e$ 的参数化。由于 $\theta$ 在 $e$ 上不分歧，$\alpha_e$ 在 $\theta_{T_1}$ 下的提升由两条不相交的道路 $\alpha^1_{e_1},\alpha^1_{e_2}\colon(0,1)\to H_{T_1}$ 组成。我们假定提升 $\alpha^1_{e_1}$ 是相对于 $H_{T_1}$ 上的谱网络从汇（sink）走向源（source）的那条提升，而 $\alpha^1_{e_2}$ 是相对于 $H_{T_1}$ 上的谱网络从源走向汇的那条提升。对 $T_2$ 我们同样有 $\alpha^2_{e_1},\alpha^2_{e_2}\colon(0,1)\to H_{T_2}$。最后，我们通过把所有 $t\in(0,1)$ 的 $\alpha^1_{e_1}(t)$ 与 $\alpha^2_{e_1}(t)$ 等同、并把 $\alpha^1_{e_2}(t)$ 与 $\alpha^2_{e_2}(t)$ 等同，来粘合六边形 $H_{T_1}$ 与 $H_{T_2}$，见图 13.11。由这个粘合构造得到的空间是一个曲面，记作 $\Sigma$。映射 $\theta_T$ 给出带谱网络的分歧覆叠 $\theta\colon\Sigma\to S$。

\Fig[0.78\linewidth]{figures/fig_13_11.png}{13.11}{沿理想三角剖分 $\Delta$ 的边 $e$ 粘合。}

% [Tutor Note: 原文把 $\theta$ 写成 $\psi_T\circ(z\mapsto z^2)\circ\phi_T^{-1}$，但按交换图表，$\theta_T$ 应把 $H_T$ 送到 $T$，正确顺序为 $\theta_T=\phi_T^{-1}\circ(z\mapsto z^2)\circ\psi_T$：先用 $\psi_T$ 到 $\HH$，再平方到 $\mathbb{T}$，最后用 $\phi_T^{-1}$ 回到 $T$。原文写作逆序，属记号疏漏；以图表为准。]

% [原文 p.152]
\subsection{高阶谱网络}

本节中，我们将利用 Goncharov 与 Kontsevich 在 [168] 中引入的秩 $k\ge2$ 丛的谱覆盖的组合构造，来给出与该 $k$ 重分歧覆叠相关联的小谱网络的组合构造。这个构造产生一类非常特殊的谱网络，在 [156] 中称为\emph{极小谱网络}；不过这一类受限的谱网络在下文仍将非常重要，特别是因为它与 Fock--Goncharov 坐标的联系（见 §15.3）。我们先给出 [168] 中谱覆盖的构造。设 $S$ 为定向纤毛曲面，赋有理想三角剖分 $\Delta$。对任何整数 $k\ge2$，我们通过对 $S$ 的所有三角形 $t$ 粘合如图 13.12 所述的基本片段 $\Gamma_k^{(t)}$，定义 $S$ 上的\emph{二分图} $\Gamma_k$。

\Fig[0.85\linewidth]{figures/fig_13_12.png}{13.12}{限制在 $\Delta$ 的三角形 $t$ 上的图 $\Gamma_2^{(t)}$、$\Gamma_3^{(t)}$ 与 $\Gamma_4^{(t)}$。}

然后对每对共享三角剖分的一条边 $e$ 的三角形 $t$ 与 $t'$，我们在 $e$ 上把 $\Gamma_k^{(t)}$ 与 $\Gamma_k^{(t')}$ 相应的黑顶点等同；见图 13.13。白顶点是三价的，而黑顶点是单价的、二价的或三价的：单价与二价的黑顶点称为\emph{外部的}，三价的黑顶点称为\emph{内部的}。单价的黑顶点恰是 $\Gamma_k$ 落在 $S$ 边界上的顶点。

\Fig[0.48\linewidth]{figures/fig_13_13.png}{13.13}{对两个相邻三角形 $t$ 与 $t'$，$\Gamma_3^{(t)}$ 与 $\Gamma_3^{(t')}$ 的粘合。}

图 $\Gamma_k$ 嵌入 $S$ 中，故每个顶点从 $S$ 的定向继承其关联边上的循环次序。设 $\gamma=(v_1,\dots,v_r)$ 为
% [原文 p.153]
$\Gamma_k$ 上的一条道路，且 $1<i<k$。若边 $(v_{i-1},v_i)$ 与 $(v_i,v_{i+1})$（相应地 $(v_i,v_{i+1})$ 与 $(v_{i-1},v_i)$）按该次序在 $v_i$ 处依循环次序相邻，则说 $\gamma$ 在 $v_i$ 处\emph{左转}（相应地\emph{右转}）。由于图 $\Gamma_k$ 嵌入 $S$，$\Gamma_k$ 中的道路将与其在 $S$ 中的像等同。

$\Gamma_k$ 上的\emph{zig-zag 路径}（zig-zag path）是一条在每处白顶点右转、在每处黑顶点左转的道路 $(v_1,\dots,v_r)$。特别地，$\Gamma_k$ 的每条定向边恰好给出一条以该边开头的 zig-zag 路径。于是每条 zig-zag 路径都可以向两侧极大延拓。由于 $\Gamma_k$ 有限，一条 zig-zag 路径或者延拓为连接两个边界分支的道路，或者延拓为 $\Gamma_k$ 上的一个圈，而该圈的任何延拓都是周期的。以下我们假定所有 zig-zag 路径都是简单闭圈或连接边界分支的道路。

% [Tutor Note: 原文「$\Gamma_k$ 上的一条道路 $\gamma=(v_1,\dots,v_r)$，且 $1<i<k$」中上界应为 $r$（路径的顶点数），即 $1<i<r$；否则「$\gamma$ 在 $v_i$ 处左转/右转」中 $i$ 的取值范围将依赖于秩 $k$ 而与路径长度无关，显然不合定义。]

\begin{proposition}[13.16]
设 $\gamma$ 是 $\Gamma_k$ 上的 zig-zag 路径。则 $\gamma$ 在 $S$ 中的像或者同胚于一个圆周 $S^1$（其定向与 $S$ 的定向相符），或者同胚于一个闭区间，且 $S\setminus\gamma$ 有两个连通分支，其中一个是同胚于带一个穿孔的开圆盘或带一个边界穿孔的开半圆盘的连通分支。我们把这个特定的连通分支记作 $S_\gamma$，并把它内部的穿孔记作 $s_\gamma$。
\end{proposition}

设 $\gamma$ 是 $\Gamma_k$ 上的 zig-zag 路径，$\bar S_\gamma$ 是 $S_\gamma$ 在 $S$ 中的闭包。$S_\gamma$ 的边界是 $\gamma$ 在 $S$ 中的像，特别地它不含任何穿孔，因此 $\bar S_\gamma$ 同胚于带一个穿孔的闭（半）圆盘。

\begin{lemma}[13.17]
设 $k\ge2$。对 $S$ 的每个穿孔 $s$，$\Gamma_k$ 中恰有 $k$ 条 zig-zag 路径 $\gamma_s^{(1)},\dots,\gamma_s^{(k)}$，使得对一切 $1\le i\le k$ 有 $s=s_{\gamma_s^{(i)}}$。此外，对一切 $1\le i<j\le k$，有 $\gamma_s^{(i)}\subset S_{\gamma_s^{(j)}}$（图 13.14）。
\end{lemma}

\Fig[0.42\linewidth]{figures/fig_13_14.png}{13.14}{$\Gamma_3^{(t)}$ 中环绕穿孔 $s$ 的三条 zig-zag 路径 $\gamma_s^{(1)}$、$\gamma_s^{(2)}$ 与 $\gamma_s^{(3)}$。$t$ 在路径 $\gamma_s^{(i)}$ 上方的部分包含在 $S_{\gamma_s^{(i)}}$ 中。}

我们定义 $\Sigma_k$ 为把带穿孔的（半）圆盘粘到 $\Gamma_k$ 上所得的拓扑曲面：沿 $\Gamma_k$ 中每条 zig-zag 路径 $\gamma$，我们粘上一个 $\bar S_\gamma$ 的副本。我们如下定义 $\pi\colon\Sigma_k\to S$：$\Gamma_k$ 中一点的像是它在 $S$ 中的浸入像，而每个胞腔 $S_\gamma$
% [原文 p.154]
赋以由包含给出的映射 $S_\gamma\to S$。我们把 $\Sigma_k$ 中位于 $S_{\gamma_s^{(i)}}$ 内部的穿孔记作 $s_i$。

\begin{proposition}[13.18]
映射 $\pi\colon\Sigma_k\to S$ 是带简单分歧点的分歧 $k$ 重覆叠。分歧点恰是 $\Gamma_k$ 的内部黑顶点。
\end{proposition}

\begin{proof}
$\Gamma_k$ 的每条边都是恰好两条 zig-zag 路径的一部分，故 $\Sigma_k$ 在 $\Gamma_k$ 每条边内部的点附近以及在外部黑顶点附近都是拓扑曲面。每个白顶点恰是三条 zig-zag 路径边界上的点，每个胞腔 $S_\gamma$ 被粘到一对相邻的边上（见图 13.15）。于是白顶点的邻域是拓扑曲面。在所有这些情形中，我们也看出覆叠 $\pi$ 是正则的。在每个内部黑顶点的邻域中，$\Sigma_k$ 也是拓扑曲面，但覆叠 $\pi$ 在每个内部黑顶点处都是简单分歧的（见图 13.15）。
\end{proof}

\Fig[0.75\linewidth]{figures/fig_13_15.png}{13.15}{白顶点、内部黑顶点与外部黑顶点的邻域，以及这些顶点邻域中粘合的不同胞腔。}

\begin{remark}[13.19]
曲面 $\Sigma_2$ 与覆叠 $\pi\colon\Sigma_2\to S$ 同胚于 §13.5.1 中构造的那些。
\end{remark}

\Fig[0.85\linewidth]{figures/fig_13_16.png}{13.16}{$S$ 的一个三角形的 $3$ 重分歧覆叠。波浪蓝线是分支割线。}

我们现在可以构造与 $k$ 重分歧覆叠 $\pi\colon\Sigma_k\to S$ 相关联的 $\Sigma_k$ 上的小谱网络。为此，我们将描述 $S$ 的一个三角形的覆叠上的谱网络，而 $\Sigma_k$ 上的完整谱网络是所有三角形覆叠上的谱网络的并。首先注意，$S$ 的每个穿孔 $s$ 都是非分歧的，即 $|\pi^{-1}(s)|=k$。$s$ 的这 $k$ 条提升 $s_1,\dots,s_k$ 各自落在
% [原文 p.155]
引理 13.17 中定义的某个带穿孔圆盘 $S_{\gamma_s^{(i)}}$ 中，并继承由这些带穿孔圆盘的包含关系给出的次序 $s_1<\dots<s_k$。设 $t$ 是 $S$ 的一个三角形，$s$ 是该三角形的一个顶点，且 $2\le i\le k$。在 $S_{\gamma_s^{(i)}}\cap t$ 的边界上有 $k-i$ 个内部黑顶点，在 $S_{\gamma_s^{(i-1)}}\cap t$ 的边界上有 $k-i$ 个白顶点。这些黑顶点每一个都通过 $\Gamma_k^{(t)}$ 的一条边恰好连接到这些白顶点之一，且这些黑顶点是分歧点。在 $\Sigma_k$ 上，我们如下为这些黑顶点各定义一条谱网络的线：设 $v_b$ 是 $S_{\gamma_s^{(i)}}\cap t$ 边界上的一个内部黑顶点，$v_w$ 是 $S_{\gamma_s^{(i-1)}}\cap t$ 边界上与 $v_b$ 关联的白顶点。设 $p_1$ 是 $S_{\gamma_s^{(i)}}\cap t$ 中从 $s_i$ 到 $v_b$ 的一条道路，$p_2$ 是由沿图 $\Gamma_k^{(t)}$ 的边得到的从 $v_b$ 到 $v_w$ 的道路，$p_3$ 是 $S_{\gamma_s^{(i-1)}}\cap t$ 中从 $v_w$ 到 $s_{i-1}$ 的道路。谱网络的线就是拼接 $p_3.p_2.p_1$（$k=3$ 的情形见图 13.16 与图 13.17）。

\begin{proposition}[13.20]
如上构造的 $\Sigma_k$ 上所有道路构成的集合 $\mathcal{W}$ 是与 $k$ 重分歧覆叠 $\pi\colon\Sigma_k\to S$ 相关联的小谱网络。
\end{proposition}

\Fig[0.85\linewidth]{figures/fig_13_17.png}{13.17}{提升到 $\Sigma_3$ 的三角形上的小谱网络。谱网络诸线依其所经过的分支点而着以不同颜色。}

\begin{proof}
每条线 $p\colon[-1,1]\to\bar\Sigma_k$ 满足 $\pi(p(t))=\pi(p(-t))$（$t\in[-1,1]$），且 $\pi(p(-1))=\pi(p(1))$ 是 $\bar S$ 中的穿孔、$p(0)$ 是内部黑顶点，即分歧点。每个内部黑顶点属于三条 zig-zag 路径，因此位于三个形如 $S_{\gamma_s^{(i)}}$ 的胞腔的边界上。过每个黑顶点有三条 $\mathcal{W}$ 的线，方向交替（见图 13.18）。
% [原文 p.156]
$\Gamma_k$ 中与内部黑顶点相接的每条边恰好被 $\mathcal{W}$ 的一条线跟随，而与外部黑顶点相接的边不与 $\mathcal{W}$ 的任何线相交。在形如 $S_{\gamma_s^{(i)}}$ 的带穿孔圆盘内，$\mathcal{W}$ 的线之间没有相交；因此 $\mathcal{W}$ 的线之间（除分歧点处的相交外）的唯一相交都位于 $\Gamma_k$ 的白顶点上。这意味着所有这些相交都发生在已经越过其分歧点的线之间，故集合 $\mathcal{W}$ 是小谱网络。
\end{proof}

\Fig[0.40\linewidth]{figures/fig_13_18.png}{13.18}{谱网络三条线在内部黑顶点处的相交。}

\section{解析构造：WKB 谱网络}

现在设带穿孔的曲面 $S$ 按注 4.10 的意义赋有复结构，并考虑其紧化 $\bar S:=S\cup P$（参看 §4.3）。设 $\pi\colon T^*\bar S\to\bar S$ 为 $\bar S$ 上的全纯余切丛，其全纯结构由 $S$ 的复结构典范地诱导。考虑一个多元组 $\boldsymbol\phi:=(\phi_1,\dots,\phi_k)$，其中 $\phi_i$ 是 $(T^*\bar S)^{\otimes i}$ 的亚纯截面（对一切 $i\in\{1,\dots,k\}$），而 $P\subset\bar S$ 是 $\boldsymbol\phi$ 的所有极点构成的集合。

与 $\boldsymbol\phi$ 相关联的\emph{谱覆盖}定义为
\[
\Sigma_{\boldsymbol\phi}:=\Bigl\{\sum_{i=1}^{k}\phi_i y^{k-i}=0\ \Big|\ y\in T^*S\Bigr\}\subset T^*S,
\]
其中 $T^*S$ 是 $S$ 的余切丛。若把上面的截面 $\phi_i$ 取得足够一般，则覆叠 $\pi|_{\Sigma_{\boldsymbol\phi}}\colon\Sigma_{\boldsymbol\phi}\to S$ 是只有简单分歧的光滑分歧 $k$ 重覆叠。如前，令 $B$ 为 $S$ 中分支点的集合。

\begin{remark}[13.21]
为简便起见，我们将滥用记号，写作 $\pi\colon\Sigma_{\boldsymbol\phi}\to S$ 而不写 $\pi|_{\Sigma_{\boldsymbol\phi}}\colon\Sigma_{\boldsymbol\phi}\to S$。
\end{remark}

\begin{definition}[13.22]
若 $\Sigma_{\boldsymbol\phi}$ 光滑且只有简单分歧点，则称覆叠 $\pi\colon\Sigma_{\boldsymbol\phi}\to S$ 是\emph{简单的}。
\end{definition}

\begin{remark}[13.23]
当 $n=2$ 且 $\phi_1=0$ 时，$\Sigma_{\boldsymbol\phi}:=\{y^2+\phi_2=0\mid y\in T^*S\}$。特别地，$\Sigma_{\boldsymbol\phi}$ 简单当且仅当 $\phi_2$ 只有单零点。
\end{remark}

% [Tutor Note: 注 13.23 中的「$n=2$」应为「$k=2$」：本章一律用 $k$ 表示覆叠的次数，$n$ 未被用作秩；且结论 $y^2+\phi_2=0$ 正是 $k=2$ 时 $\sum_i\phi_i y^{k-i}=0$ 的形式。]

设如前 $\pi\colon\Sigma\to S$ 是只有简单分歧的光滑谱覆盖。我们定义空间
\[
\hat\Sigma:=\bigl\{(y,y')\in\Sigma\times\Sigma\ \big|\ \pi(y)=\pi(y'),\ y\neq y'\bigr\}\subset T^*S\times T^*S,
\]
称为覆叠 $\pi\colon\Sigma\to S$ 的\emph{根覆盖}（root covering）。它是 $S$ 的光滑 $(k^2-k)$ 叶分歧覆叠。存在自然的对合 $\mu\colon\hat\Sigma\to\hat\Sigma$，定义为 $\mu(y,y'):=(y',y)$，以及两个投影映射 $p_1,p_2\colon\hat\Sigma\to\Sigma$，分别满足 $p_1(x,y):=x$ 与 $p_2(x,y):=y$。显然 $p_1=p_2\circ\mu$ 且 $\pi\circ p_1=\pi\circ p_2:=\hat\pi$。

% [原文 p.157]
在根覆盖上存在自然的阿贝尔微分 $\rho$，定义如下：设 $\lambda$ 是 $T^*S$ 上的典范刘维尔（Liouville）$1$-形式。我们定义 $\rho:=p_1^*\lambda-p_2^*\lambda$，它在分歧点上方有零点。若取对应于 $\hat\Sigma$ 的定向叶状结构 $F_\rho$，则 $\mu$ 保持 $F_\rho$ 的叶，但反转定向，因为 $\mu^*\rho=-\rho$。

我们按以下两步，把与多元组 $\boldsymbol\phi$ 相关联的谱网络 $\mathcal{W}(\boldsymbol\phi)$ 构造为 $F_\rho$ 的一族叶：
\begin{enumerate}
\item $\rho$ 的临界图包含在 $\mathcal{W}(\boldsymbol\phi)$ 中；
\item 设 $\alpha_1,\alpha_2\colon(-1,1)\to\hat\Sigma$ 是 $\mathcal{W}(\boldsymbol\phi)$ 的两片叶，使得存在 $t_1,t_2\in(-1,1)$ 满足：$\alpha_1(t_1)$ 位于叶（sheet）$(i,j)$ 中，$\alpha_2(t_2)$ 位于叶 $(j,k)$ 中，且 $p_j(\alpha_1(t_1))=p_j(\alpha_2(t_2))$，即这两片叶到 $\Sigma$ 的投影在叶 $j$ 中相交。设 $p_0$ 为叶 $(i,k)$ 中的点，$p'_0$ 为叶 $(k,i)$ 中的点，使得 $\hat\pi(p_0)=\hat\pi(p'_0)=\hat\pi(\alpha_1(t_1))\in S$。然后我们向 $\mathcal{W}(\boldsymbol\phi)$ 添加 $F_\rho$ 的两片半叶：一片在叶 $(i,k)$ 中、从点 $p_0$ 出发且以 $p_0$ 为该叶的源；另一片在叶 $(k,i)$ 中、从点 $p'_0$ 出发且以 $p'_0$ 为该叶的汇；见图 13.19。
\end{enumerate}

\Fig[0.60\linewidth]{figures/fig_13_19.png}{13.19}{上文第 2 步所述的构造。}

% [原文 p.158]
上述构造 $\mathcal{W}(\boldsymbol\phi)$ 的过程一般地是无穷的，因为第二步中添加一片叶通常会产生新的叶在 $\Sigma$ 中投影后的相交。但有时该过程在有限多步后终止。在后一情形，$\mathcal{W}(\boldsymbol\phi)$ 投影到 $\Sigma$ 与 $S$ 后满足谱网络的条件。我们把对应于亚纯微分的谱网络称为 \emph{WKB 谱网络}。

\begin{example}[13.24]
若 $k=2$，则 $p_1$ 与 $p_2$ 是曲面的微分同胚。

若 $k=3$，则 $p_1\circ\pi\colon\hat\Sigma\to X$ 是 $6:1$ 分歧覆叠，$\pi\colon\Sigma\to X$ 是 $3:1$ 分歧覆叠，而 $p_i\colon\hat\Sigma\to\Sigma$（$i\in\{1,2\}$）是 $2:1$ 分歧覆叠。
\end{example}

% [Tutor Note: 例 13.24 中两处把底曲面写成 $X$（$p_1\circ\pi\colon\hat\Sigma\to X$、$\pi\colon\Sigma\to X$），而本章其余各处均记底曲面为 $S$；此处应为 $S$。它不影响三个次数的计算：$k=3$ 时 $|\hat\Sigma\to S|=k(k-1)=6$，$|\Sigma\to S|=3$，$|\hat\Sigma\to\Sigma|=2$。]

\begin{example}[13.25]
设 $\bar S=\C\PP^1$ 且 $k=3$。取 $\boldsymbol\phi=\ee^{\frac{i\pi}{2}}(0,-\dd z^2,z\,\dd z^3)$。$\boldsymbol\phi$ 的唯一极点是无穷远点，故 $S=\C\PP^1\setminus\{\infty\}=\C$。于是
\[
\Sigma=\bigl\{(z,y)\in T^*S\ \big|\ y^3-y+z=0\bigr\}.
\]
分支轨迹与判别式为零的点集一致：$\Delta(z)=4z^2-27=0$，由两个点组成。谱网络 $\mathcal{W}(\boldsymbol\phi)$ 就是图 13.9 中所示的那一个。
\end{example}

% [Tutor Note: 对 $y^3-y+z$，按判别式公式 $\Delta=-4a^3-27b^2$（$a=-1$，$b=z$）得 $\Delta(z)=4-27z^2$，零点为 $z=\pm2/(3\sqrt3)$。原文写 $\Delta(z)=4z^2-27$，相当于把系数次序写反（两种写法都恰有两个零点，故「由两个点组成」无误）。核对分支点位置时以 $4-27z^2$ 为准。]

\begin{remark}[13.26]
谱网络的更多例子见第 22 章，那里也讨论了 WKB 谱网络的物理实质。
\end{remark}

\begin{remark}[13.27]
所有秩 $2$ 的谱网络都同痕于一个 WKB 谱网络。反之，由于纤毛曲面上的二次微分定义一个理想三角剖分（见 §8.3），任何秩 $2$ 的 WKB 谱网络都具有 §13.5.1 中所述的形式。相反，对秩 $k\ge3$ 的谱网络，这两个论断都不成立：并非所有非退化谱网络都同痕于 WKB 谱网络，也并非所有 WKB 谱网络都具有 §13.5.2 中所述的形式。
\end{remark}

\section{用非退化谱网络作道路提升}

现在既然有了小谱网络的定义，我们就可以描述道路提升规则，它使后文能够构造非阿贝尔化映射。

% [原文 p.159]
设 $\pi\colon\Sigma\to S$ 为满足上述条件的分歧 $k$ 重覆叠，$\mathcal{W}$ 为 $\Sigma$ 上秩 $k$ 的小谱网络。我们希望构造一个映射 $\pi_1(S,p)\to\Z^{\mathrm{tw}}_B[\mathcal{P}(\Sigma)]$，其中 $p\in S$ 且 $p$ 不在 $\mathcal{W}$ 的任何线上。

我们考虑道路 $\gamma\colon[0,1]\to S$，其端点不在谱网络的任何线上，也不经过任何分支点。为把这条道路提升到 $\bar\Sigma$，我们把它分成若干片段 $\gamma=\gamma_n\cdot\ldots\cdot\gamma_1$，使每个 $\gamma_i$ 至多与谱网络 $\mathcal{W}$ 相交一次。这里我们用 $\cdot$ 表示道路的拼接，如路径代数的定义中那样。

若 $\gamma_i$ 不与谱网络相交，则按通常方式提升，得到 $\Sigma$ 中同样不与谱网络相交的 $k$ 条提升 $\gamma_i^{(1)},\dots,\gamma_i^{(k)}$。

若 $\gamma_i$ 与谱网络的线 $p_h$ 相交，则 $\gamma_i$ 只有两条标准提升在 $\Sigma$ 中与 $p_h$ 相交。这两条提升中的一条 $\gamma_i^{(j)}$ 在 $p_h$ 经过其分歧点\emph{之前}与 $p_h$ 相交，另一条 $\gamma_i^{(\ell)}$ 则在 $p_h$ 经过其分歧点\emph{之后}与之相交。在这种情形下，我们必须在标准提升之外再加一条新道路 $\gamma'_i$，如图 13.20。道路 $\gamma'_i$ 这样得到：沿 $\gamma_i^{(j)}$ 从其起点走到它与线 $p_h$ 的交点，然后沿 $p_h$（不穿过它）走到它与 $\gamma_i^{(\ell)}$ 相遇之处，再沿 $\gamma_i^{(\ell)}$ 走到其终点。

\Fig[0.65\linewidth]{figures/fig_13_20.png}{13.20}{用谱网络作道路提升。图为 $\Sigma$ 中（左）与 $S$ 中（右）的情形。}

\begin{definition}[13.28]
曲线 $\gamma$ 在 $\Sigma$ 中相对于谱网络 $\mathcal{W}$ 的\emph{提升} $\tilde\gamma$，是 $\gamma_i$ 的所有提升在 $\Z^{\mathrm{tw}}_B[\mathcal{P}(\Sigma)]$ 中的乘积：
\[
\tilde\gamma=\Bigl(\sum_{i=1}^{k}\gamma_1^{(i)}+\gamma'_1\Bigr)\cdot\ldots\cdot\Bigl(\sum_{i=1}^{k}\gamma_n^{(i)}+\gamma'_n\Bigr),
\]
其中对不与谱网络相交的片段 $\gamma_i$，令 $\gamma'_i=0$。
\end{definition}

% [Tutor Note: 定义 13.28 中求和号只作用于 $\gamma_1^{(i)}$，$\gamma'_1$ 在求和号外；若该片段不与 $\mathcal{W}$ 相交则取 $\gamma'_1=0$，此时括号化为 $\sum_{i=1}^k\gamma_1^{(i)}$。核对公式时注意 $n$ 是片段数、$k$ 是秩，两者不可混。]

\begin{theorem}[13.29]
这条道路提升规则是同伦不变的，即若两条道路 $\gamma$ 与 $\delta$ 相对于其端点同伦，则在 $\Z^{\mathrm{tw}}_B[\mathcal{P}(\Sigma)]$ 中 $\tilde\gamma=\tilde\delta$。
\end{theorem}

该定理是两个下述引理的推论：

\begin{lemma}[13.30]
设 $\gamma$ 是 $S$ 上的道路，它恰与谱网络的同一条线 $p_h$ 相交两次、且不与 $\mathcal{W}$ 的其他线相交，如图 13.21。则 $\tilde\gamma=\gamma^{(1)}+\dots+\gamma^{(k)}$，其中 $\gamma^{(1)},\dots,\gamma^{(k)}$ 是 $\gamma$ 到 $\Sigma$ 的 $k$ 条标准提升。
\end{lemma}

\begin{proof}
谱网络产生了另外两条提升 $\gamma'$ 与 $\gamma''$。这两条特殊提升在越过一个分歧点的意义下同伦，故在 $\Z^{\mathrm{tw}}_B[\mathcal{P}(\Sigma)]$ 中 $\gamma'+\gamma''=0$；见图 13.22。
\end{proof}

% [Tutor Note: $\gamma'+\gamma''=0$ 的机制：这两条额外提升只相差绕过分歧点的一个小环路 $\alpha_z^b$，而在扭转路径代数中由定义有 $\alpha_z^b=-e_z$，故 $\gamma'+\gamma''=0$；引理 13.31 证明中的 $\gamma'_1+\gamma'_5=0$ 等抵消同样归结为关系 $\alpha_z^b+e_z=0$。]

% [原文 p.160]
\Fig[0.40\linewidth]{figures/fig_13_21.png}{13.21}{与 $\mathcal{W}$ 的同一根线相交两次的环路。}

\Fig[0.72\linewidth]{figures/fig_13_22.png}{13.22}{$\gamma$ 的谱网络提升。}

\begin{lemma}[13.31]
设 $\gamma$ 是以 $m\in S$ 为基点的可缩环路，它绕 $S$ 中的分支点 $b$ 一圈，且恰与 $\mathcal{W}$ 从 $b$ 出发的三条线各相交一次，如图 13.23。则 $\tilde\gamma=e_{m_1}+\dots+e_{m_k}$，其中 $m_1,\dots,m_k$ 是 $m$ 到 $\Sigma$ 的 $k$ 条提升，而 $e_{m_i}$ 表示以 $m_i$ 为基点的平凡环路。
\end{lemma}

\Fig[0.40\linewidth]{figures/fig_13_23.png}{13.23}{绕分支点的小环路。}

\begin{proof}
对 $\gamma$ 应用谱网络提升规则，得到 $k+6$ 条道路：$k$ 条标准提升 $\gamma^{(1)},\dots,\gamma^{(k)}$（其中 $\gamma^{(j)}$ 与 $\gamma^{(\ell)}$ 如前与 $\mathcal{W}$ 相交），以及图 13.24 中所示的六条额外道路 $\gamma'_1,\dots,\gamma'_6$。在重新标号的意义下，我们假定 $\gamma^{(j)}$ 与 $\gamma^{(\ell)}$ 的端点为 $m_1$ 与 $m_2$。在 $\Z^{\mathrm{tw}}_B[\mathcal{P}(\Sigma)]$ 中有：
% [原文 p.161]
\[
\begin{aligned}
&\gamma^{(\ell)}+\gamma'_3=0,\\
&\gamma^{(j)}+\gamma'_6=0,\\
&\gamma'_1+\gamma'_5=0,\\
&\gamma'_4=e_{m_1},\\
&\gamma'_2=e_{m_2}.
\end{aligned}
\]
因此我们可以写
\[
\tilde\gamma=\sum_{i=1}^{k}\gamma^{(i)}+\gamma'_1+\gamma'_2+\gamma'_3+\gamma'_4+\gamma'_5+\gamma'_6=\sum_{i=1}^{k}e_{m_i}.
\]
\end{proof}

\Fig[0.68\linewidth]{figures/fig_13_24.png}{13.24}{由与谱网络相交而添加的全部六条道路，以及标准提升 $\gamma^{(j)}$ 与 $\gamma^{(\ell)}$。左上图中是与平凡道路同伦的道路，其余各图中是剩下的提升，它们按在 $\Z^{\mathrm{tw}}_B[\mathcal{P}(\Sigma)]$ 中相互抵消的道路对分组。}

% [原文 p.162]
\begin{exercise}
道路提升规则可以推广到非退化谱网络，并且它是同伦不变的。
\end{exercise}

我们给出如何提升一条与线 $p_h$（$h\in H_2$）相交的道路的提示。在图 13.25 中，我们用红色画出 $S$ 中正在被提升的曲线及它们在 $\Sigma$ 中的标准提升，并用蓝色画出额外的提升。

\Fig[0.78\linewidth]{figures/fig_13_25.png}{13.25}{沿关节的道路提升。}
